\subsection{Matrix Representations}
\begin{definition} \label{def:matrices-and-matrix-operations} \hlt{Matrices and Matrix Operations}
\begin{enumerate}[label=(\roman*)]
\setlength{\itemsep}{0pt}
\item \hlt{Matrix space}: for $m,n\in\mathbb{Z}_{>0}$, $\operatorname{Mat}(m,n,\mathbb{F})=\{[A_{j,k}]_{1\leq j\leq m,1\leq k\leq n}:A_{j,k}\in\mathbb{F}\}$.
\item \hlt{Entries, rows, columns}: for $A=[A_{j,k}]\in\operatorname{Mat}(m,n,\mathbb{F})$,\\
$A_{j,k}$ is the $(j,k)$-entry, $A_{j,\cdot}$ is row $j$, and $A_{\cdot,k}$ is column $k$.
\item \hlt{Transpose}: for $A\in\operatorname{Mat}(m,n,\mathbb{F})$, $A^t\in\operatorname{Mat}(n,m,\mathbb{F})$ is defined by $(A^t)_{k,j}=A_{j,k}$.
\item \hlt{Main diagonal}: for $A=[A_{j,k}]\in\operatorname{Mat}(m,n,\mathbb{F})$ and $r=\min(m,n)$, $\operatorname{diag}(A)=(A_{1,1},\ldots,A_{r,r})$.
\item \hlt{Symmetry}: for $A\in\operatorname{Mat}(n,n,\mathbb{F})$, $A$ is symmetric if $A^t=A$ and skew-symmetric if $A^t=-A$.
\item \hlt{Triangularity}: for $A\in\operatorname{Mat}(n,n,\mathbb{F})$, $A$ is upper triangular if $A_{j,k}=0$ for $j>k$,\\
lower triangular if $A_{j,k}=0$ for $j<k$, and diagonal if $A_{j,k}=0$ for $j\neq k$.
\item \hlt{Linear operations}: for $A,C\in\operatorname{Mat}(m,n,\mathbb{F})$ and $a\in\mathbb{F}$, $(A+C)_{j,k}=A_{j,k}+C_{j,k}$ and $(aA)_{j,k}=aA_{j,k}$.
\item \hlt{Matrix product}: for $A\in\operatorname{Mat}(m,n,\mathbb{F})$, $B\in\operatorname{Mat}(n,p,\mathbb{F})$, $AB\in\operatorname{Mat}(m,p,\mathbb{F})$ is $(AB)_{j,k}\defeq\sum\limits_{r=1}^{n}A_{j,r}B_{r,k}$.
\item \hlt{Identity and inverse}: $I_n=[\delta_{j,k}]\in\operatorname{Mat}(n,n,\mathbb{F})$; for $A,B\in\operatorname{Mat}(n,n,\mathbb{F})$, $B$ is a two-sided inverse of $A$ if $AB=I_n=BA$, and $A$ is invertible if such $B$ exists.
\end{enumerate}
\end{definition}

\begin{lemma} \label{lem:two-sided-inverse-uniqueness} \hlt{Uniqueness of Two-Sided Inverses}\\
Let $A,B,C\in\operatorname{Mat}(n,n,\mathbb{F})$. If $AB=I_n$ and $BC=I_n$, then $A=C$.
\begin{proof}
Using $BC=I_n$ and $AB=I_n$, $A=AI_n=A(BC)=(AB)C=I_nC=C$.\\
\hlt{Structural Commentary}: A common middle inverse forces the left and right inverses to coincide.\\
\hlt{Key Insight}: Associativity turns $AB=I_n=BC$ into $A=C$.
\end{proof}
\end{lemma}

\begin{definition} \hlt{Inverse Matrix Notation}\\
If $A\in\operatorname{Mat}(n,n,\mathbb{F})$ is invertible, Lemma \ref{lem:two-sided-inverse-uniqueness} shows that its two-sided inverse is unique.\\
This inverse is denoted $A^{-1}$, and the set of invertible matrices in $\operatorname{Mat}(n,n,\mathbb{F})$ is denoted $\operatorname{Mat}(n,n,\mathbb{F})^{\times}$.
\end{definition}

\begin{proposition} \label{prop:matrix-spaces-are-vector-spaces} \hlt{Matrix Spaces}\\
For positive integers $m,n$, $\operatorname{Mat}(m,n,\mathbb{F})$ is a vector space over $\mathbb{F}$ with dimension $mn$.
\begin{proof}
Entrywise addition and scalar multiplication satisfy the vector space axioms because $\mathbb{F}$ does. For $1\leq j\leq m$ and $1\leq k\leq n$, let $E_{j,k}$ be the matrix with $1$ in entry $(j,k)$ and $0$ elsewhere.\\
Every $A=[A_{j,k}]$ has the unique expansion $A=\sum_{j=1}^m\sum\limits_{k=1}^n A_{j,k}E_{j,k}$.\\
Thus the $mn$ matrices $E_{j,k}$ form a basis, so $\dim(\operatorname{Mat}(m,n,\mathbb{F}))=mn$.\\
\hlt{Structural Commentary}: A matrix is a finite array of independent scalar slots.\\
\hlt{Key Insight}: The standard entry matrices isolate those slots one at a time.
\end{proof}
\end{proposition}

\begin{proposition} \label{prop:matrix-product-entry-row-column} \hlt{Matrix Product Entries Are Row Times Column}\\
If $A\in\operatorname{Mat}(m,n,\mathbb{F})$ and $B\in\operatorname{Mat}(n,p,\mathbb{F})$, then $(AB)_{j,k}=A_{j,\cdot}B_{\cdot,k}$ for $1\leq j\leq m$ and $1\leq k\leq p$.
\begin{proof}
The row-column product $A_{j,\cdot}B_{\cdot,k}$ is exactly $\sum_{r=1}^n A_{j,r}B_{r,k}$, which is the definition of $(AB)_{j,k}$.\\
\hlt{Structural Commentary}: Matrix multiplication is assembled entry-by-entry from row-column pairings.\\
\hlt{Key Insight}: Each entry of $AB$ sees one row of $A$ and one column of $B$.
\end{proof}
\end{proposition}

\begin{proposition} \label{prop:matrix-vector-product-column-combination} \hlt{Matrix-Vector Products as Column Combinations}\\
If $A=[\mathbf{a}_1\ \cdots\ \mathbf{a}_n]\in\operatorname{Mat}(m,n,\mathbb{F})$ and $\mathbf{x}=(x_1,\ldots,x_n)^t\in\mathbb{F}^n$, then $A\mathbf{x}=x_1\mathbf{a}_1+\cdots+x_n\mathbf{a}_n$.
\begin{proof}
The $j$th entry of $A\mathbf{x}$ is $\sum_{k=1}^n A_{j,k}x_k$, which is also the $j$th entry of $x_1\mathbf{a}_1+\cdots+x_n\mathbf{a}_n$.\\
\hlt{Structural Commentary}: Multiplying by a column selects a linear combination of the columns of $A$.\\
\hlt{Key Insight}: The coordinates of $\mathbf{x}$ are precisely the column coefficients.
\end{proof}
\end{proposition}

\begin{definition} \label{def:column-space-matrix-rank-nullity} \hlt{Column Space, Matrix Rank, and Nullity}\\
Let $A=[\mathbf{a}_1\ \cdots\ \mathbf{a}_n]\in\operatorname{Mat}(m,n,\mathbb{F})$.
\begin{enumerate}[label=(\roman*)]
\setlength{\itemsep}{0pt}
\item $\operatorname{col}(A)\defeq\operatorname{span}(\mathbf{a}_1,\ldots,\mathbf{a}_n)\subseteq\mathbb{F}^m$, and $\operatorname{row}(A)\defeq\operatorname{col}(A^t)\subseteq\mathbb{F}^n$.
\item $\operatorname{rank}(A)\defeq\dim(\operatorname{col}(A))$; column rank is $\dim(\operatorname{col}(A))$, and row rank is $\dim(\operatorname{row}(A))$.
\item $A$ has \hlt{full rank} if $\operatorname{rank}(A)=\min(m,n)$, \hlt{full column rank} if $\operatorname{rank}(A)=n$, and \hlt{full row rank} if $\operatorname{rank}(A)=m$.
\item $\ker(A)\defeq\{\mathbf{x}\in\mathbb{F}^n:A\mathbf{x}=\mathbf{0}\}$ is the \hlt{nullspace}, and $\operatorname{nullity}(A)\defeq\dim(\ker(A))$.
\end{enumerate}
\end{definition}

\begin{proposition} \label{prop:linear-system-solvability-and-affine-solution-set} \hlt{Solvability and Complete Solution Sets}\\
Let $A=[\mathbf{a}_1\ \cdots\ \mathbf{a}_n]\in\operatorname{Mat}(m,n,\mathbb{F})$ and let $\mathbf{b}\in\mathbb{F}^m$. Then $A\mathbf{x}=\mathbf{b}$ is solvable if and only if $\mathbf{b}\in\operatorname{col}(A)$.\\
If $\mathbf{x}_p$ is one solution, then every solution is of the form $\mathbf{x}=\mathbf{x}_p+\mathbf{x}_0$, where $\mathbf{x}_0\in\ker(A)$.
\begin{proof}
\begin{enumerate}[label=\roman*.]
\setlength{\itemsep}{0pt}
\item[$(\Rightarrow)$] If $A\mathbf{x}=\mathbf{b}$ for some $\mathbf{x}=(x_1,\ldots,x_n)^T$, then $\mathbf{b}=A\mathbf{x}=\sum\limits_{j=1}^n x_j\mathbf{a}_j\in\operatorname{col}(A)$.
\item[$(\Leftarrow)$] If $\mathbf{b}\in\operatorname{col}(A)$, then $\mathbf{b}=\sum_{j=1}^n x_j\mathbf{a}_j$ for some scalars $x_j$. Hence, for $\mathbf{x}=(x_1,\ldots,x_n)^T$, we have $A\mathbf{x}=\mathbf{b}$.
\end{enumerate}
Now suppose $A\mathbf{x}_p=\mathbf{b}$. For any $\mathbf{x}$, $A\mathbf{x}=\mathbf{b}\iff A(\mathbf{x}-\mathbf{x}_p)=\mathbf{0}\iff \mathbf{x}-\mathbf{x}_p\in\ker(A)$.\\
Thus every solution is exactly of the form $\mathbf{x}=\mathbf{x}_p+\mathbf{x}_0$, where $\mathbf{x}_0\in\ker(A)$.\\
\hlt{Structural Commentary}: A consistent system is an affine shift of $\ker(A)$ by any particular solution.\\
\hlt{Key Insight}: Column space decides existence; kernel decides nonuniqueness.
\end{proof}
\end{proposition}

\begin{proposition} \label{prop:matrix-product-by-columns} \hlt{Matrix Products by Columns}\\
If $A\in\operatorname{Mat}(m,r,\mathbb{F})$ and $B=[\mathbf{b}_1\ \cdots\ \mathbf{b}_n]\in\operatorname{Mat}(r,n,\mathbb{F})$, then $AB=[A\mathbf{b}_1\ \cdots\ A\mathbf{b}_n]$.
\begin{proof}
For each $k$, the $k$th column of $AB$ has entries $\sum_{\ell=1}^r A_{j,\ell}B_{\ell,k}$, which are exactly the entries of $A\mathbf{b}_k$.\\
\hlt{Structural Commentary}: Matrix products apply $A$ to columns of $B$.\\
\hlt{Key Insight}: The columns of $AB$ are obtained by applying $A$ to the columns of $B$.
\end{proof}
\end{proposition}

\begin{proposition} \label{prop:matrix-product-outer-product-decomposition} \hlt{Matrix Products as Sums of Outer Products}\\
Let $A=[\mathbf{a}_1\ \cdots\ \mathbf{a}_r]\in\operatorname{Mat}(m,r,\mathbb{F})$, and write $B\in\operatorname{Mat}(r,n,\mathbb{F})$ by rows as
\begin{align}
B=
\begin{bmatrix}
\mathbf{b}_1^t\\
\vdots\\
\mathbf{b}_r^t
\end{bmatrix}.
\nonumber
\end{align}
Then $AB=\mathbf{a}_1\mathbf{b}_1^t+\cdots+\mathbf{a}_r\mathbf{b}_r^t$.
\begin{proof}
The $(j,k)$ entry of $\mathbf{a}_\ell\mathbf{b}_\ell^t$ is $A_{j,\ell}B_{\ell,k}$. Summing over $\ell=1,\ldots,r$ gives $\sum_{\ell=1}^r A_{j,\ell}B_{\ell,k}=(AB)_{j,k}$.\\
\hlt{Structural Commentary}: Column-row products decompose a matrix product into single column-row pieces.\\
\hlt{Key Insight}: Each shared index $\ell$ contributes one outer product.
\end{proof}
\end{proposition}

\begin{proposition} \label{prop:elementary-transpose-rules} \hlt{Elementary Transpose Rules}\\
Let $A,C\in\operatorname{Mat}(m,n,\mathbb{F})$ and let $a\in\mathbb{F}$. Then $(A^t)^t=A$; $(A+C)^t=A^t+C^t$; and $(aA)^t=aA^t$.
\begin{proof}
For all admissible indices, $((A^t)^t)_{j,k}=(A^t)_{k,j}=A_{j,k}$, so $(A^t)^t=A$.\\
Also $((A+C)^t)_{j,k}=(A+C)_{k,j}=A_{k,j}+C_{k,j}=(A^t+C^t)_{j,k}$ and $((aA)^t)_{j,k}=(aA)_{k,j}=aA_{k,j}=(aA^t)_{j,k}$.\\
\hlt{Structural Commentary}: Transpose swaps indices; entrywise algebra is unchanged.\\
\hlt{Key Insight}: Transpose is linear and involutive.
\end{proof}
\end{proposition}

\begin{example} \label{ex:matrix-space-symmetric-skew-direct-sum} \hlt{Symmetric and Skew-Symmetric Decomposition}\\
Assume $\operatorname{char}(\mathbb{F})\neq2$, $M=\operatorname{Mat}(n,n,\mathbb{F})$, let $\operatorname{Sym}_n(\mathbb{F})=\{A\in M:A^t=A\}$, $\operatorname{Skew}_n(\mathbb{F})=\{A\in M:A^t=-A\}$.\\
For the map $\tau:M\to M$ defined by $\tau(A)=A^t$, linearity follows from Proposition \ref{prop:elementary-transpose-rules}. Hence $\operatorname{Sym}_n(\mathbb{F})=\ker(\tau-I)$ and $\operatorname{Skew}_n(\mathbb{F})=\ker(\tau+I)$, so both are subspaces of $M$.\\
For every $A\in M$, $A=\frac12(A+A^t)+\frac12(A-A^t) \in\operatorname{Sym}_n(\mathbb{F})+\operatorname{Skew}_n(\mathbb{F})$.\\
Also $A\in\operatorname{Sym}_n(\mathbb{F})\cap\operatorname{Skew}_n(\mathbb{F})$ implies $A=A^t=-A$, hence $A=0$. Thus $M=\operatorname{Sym}_n(\mathbb{F})\oplus\operatorname{Skew}_n(\mathbb{F})$.\\
\hlt{Structural Commentary}: Transpose splits a matrix into its even and odd parts.\\
\hlt{Key Insight}: The directness of the split uses the ability to divide by $2$.
\end{example}

\begin{proposition} \label{prop:transpose-reverses-products} \hlt{Transpose Reverses Products}\\
If $A\in\operatorname{Mat}(m,n,\mathbb{F})$ and $B\in\operatorname{Mat}(n,p,\mathbb{F})$, then $(AB)^t=B^tA^t$.
\begin{proof}
For $1\leq j\leq p$ and $1\leq k\leq m$, $((AB)^t)_{j,k}=(AB)_{k,j}=\sum\limits_{\ell=1}^n A_{k,\ell}B_{\ell,j}=\sum\limits_{\ell=1}^n (B^t)_{j,\ell}(A^t)_{\ell,k}=(B^tA^t)_{j,k}$.\\
\hlt{Structural Commentary}: Transposition swaps the order of composition encoded by matrix multiplication.\\
\hlt{Key Insight}: Rows become columns, so the product order reverses.
\end{proof}
\end{proposition}

\begin{corollary} \label{cor:matrix-inverse-product-and-transpose} \hlt{Matrix Inverses and Transposes Reverse Products}\\
Let $A,B\in\operatorname{Mat}(n,n,\mathbb{F})$ be invertible. Then $AB$, $A^t$ are invertible, and $(AB)^{-1}=B^{-1}A^{-1}$; $(A^t)^{-1}=(A^{-1})^t$.
\begin{proof}
First, $(B^{-1}A^{-1})(AB)=B^{-1}(A^{-1}A)B=B^{-1}B=I_n$; $(AB)(B^{-1}A^{-1})=A(BB^{-1})A^{-1}=AA^{-1}=I_n$.\\
Hence $(AB)^{-1}=B^{-1}A^{-1}$. By Proposition \ref{prop:transpose-reverses-products}, $(A^{-1})^tA^t=(AA^{-1})^t=I_n$; $A^t(A^{-1})^t=(A^{-1}A)^t=I_n$.\\
\hlt{Structural Commentary}: Inversion and transposition both reverse multiplication order.\\
\hlt{Key Insight}: Apply the defining inverse equations and reverse products.
\end{proof}
\end{corollary}

\begin{definition} \hlt{Coordinate Vector}\\
Let $\beta=(\mathbf{v}_1,\ldots,\mathbf{v}_n)$ be a basis of a finite-dimensional vector space $\mathcal{V}$.\\
If $\mathbf{v}=c_1\mathbf{v}_1+\cdots+c_n\mathbf{v}_n$, then the \hlt{$\beta$-coordinate vector} of $\mathbf{v}$ is
\begin{align}
[\mathbf{v}]_{\beta}\defeq
\begin{bmatrix}
c_1\\
\vdots\\
c_n
\end{bmatrix}
\in\mathbb{F}^n. \nonumber
\end{align}
The function $[\cdot]_{\beta}:\mathcal{V}\to\mathbb{F}^n$ is the \hlt{$\beta$-basis representation function}.
\end{definition}

\begin{proposition} \label{prop:coordinate-map-isomorphism} \hlt{Coordinate Maps Are Linear Bijections}\\
Let $\beta=(\mathbf{v}_1,\ldots,\mathbf{v}_n)$ be a basis of a finite-dimensional vector space $\mathcal{V}$.\\
Then the coordinate map $[\cdot]_{\beta}:\mathcal{V}\to\mathbb{F}^n$ is linear and bijective.
\begin{proof}
If $\mathbf{u}=a_1\mathbf{v}_1+\cdots+a_n\mathbf{v}_n$ and $\mathbf{v}=b_1\mathbf{v}_1+\cdots+b_n\mathbf{v}_n$, then for $c\in\mathbb{F}$,
\begin{align}
[c\mathbf{u}+\mathbf{v}]_{\beta} =
\begin{bmatrix}
ca_1+b_1\\
\vdots\\
ca_n+b_n
\end{bmatrix}
= c[\mathbf{u}]_{\beta}+[\mathbf{v}]_{\beta} \nonumber
\end{align}
Thus $[\cdot]_{\beta}$ is linear. It is injective because $[\mathbf{v}]_{\beta}=\mathbf{0}$ forces all basis coefficients of $\mathbf{v}$ to be zero, and it is surjective because every column $(c_1,\ldots,c_n)^t\in\mathbb{F}^n$ equals $[c_1\mathbf{v}_1+\cdots+c_n\mathbf{v}_n]_{\beta}$.\\
\hlt{Structural Commentary}: A basis converts abstract vectors into coordinate columns without losing structure.\\
\hlt{Key Insight}: Coordinates give a structure-preserving copy of the vector space, not extra information.
\end{proof}
\end{proposition}

\begin{definition} \hlt{Matrix Representation}\\
Let $\beta = (\mathbf{v}_1, \ldots, \mathbf{v}_n)$ be a basis of $\mathcal{V}$ and $\gamma = (\mathbf{w}_1, \ldots, \mathbf{w}_m)$ a basis of $\mathcal{W}$. For $T \in \LL(\mathcal{V}, \mathcal{W})$, the matrix of $T$ relative to $(\beta,\gamma)$ is $[T]_{\gamma \leftarrow \beta} \defeq
\begin{bmatrix}
[T(\mathbf{v}_1)]_{\gamma} & \cdots & [T(\mathbf{v}_n)]_{\gamma}
\end{bmatrix}
\in \operatorname{Mat}(m,n,\mathbb{F})$.
\end{definition}

\begin{proposition} \label{prop:stacked-nullspaces-and-concatenated-column-spaces} \hlt{Stacks Intersect Kernels, Concatenations Sum Columns}\\
Let $A\in\operatorname{Mat}(m,n,\mathbb{F})$, $B\in\operatorname{Mat}(p,n,\mathbb{F})$, $C\in\operatorname{Mat}(m,r,\mathbb{F})$, and $D\in\operatorname{Mat}(m,s,\mathbb{F})$.\\
Then $\ker\begin{bmatrix} A\\ B\end{bmatrix}=\ker(A)\cap\ker(B)$; and $\operatorname{col}\begin{bmatrix} C & D\end{bmatrix} =\operatorname{col}(C)+\operatorname{col}(D)$.
\begin{proof}
For $\mathbf{x}\in\mathbb{F}^n$, $\mathbf{x}\in\ker\begin{bmatrix}A\\ B\end{bmatrix} \iff \begin{bmatrix}A\\ B\end{bmatrix}\mathbf{x}=\mathbf{0} \iff A\mathbf{x}=\mathbf{0}\ \text{and}\ B\mathbf{x}=\mathbf{0} \iff \mathbf{x}\in\ker(A)\cap\ker(B)$.\\
For $\mathbf{y}\in\mathbb{F}^m$, $\mathbf{y}\in\operatorname{col}\begin{bmatrix}C&D\end{bmatrix} \iff \exists\,\mathbf{u}\in\mathbb{F}^r,\ \mathbf{v}\in\mathbb{F}^s$ such that $\mathbf{y}=\begin{bmatrix}C&D\end{bmatrix}\begin{bmatrix}\mathbf{u}\\ \mathbf{v}\end{bmatrix}$\\
$\iff \exists\,\mathbf{u},\mathbf{v}$ such that $\mathbf{y}=C\mathbf{u}+D\mathbf{v} \iff \mathbf{y}\in\operatorname{col}(C)+\operatorname{col}(D)$.\\
\hlt{Structural Commentary}: Stacking intersects constraints; concatenation adds column directions.\\
\hlt{Key Insight}: Stacks encode intersections of nullspaces, while concatenations encode sums of column spaces.
\end{proof}
\end{proposition}

\begin{proposition} \label{prop:matrix-representation-respects-linear-combinations} \hlt{Matrices Respect Linear Combinations of Maps}\\
Let $\beta=(\mathbf{v}_1,\ldots,\mathbf{v}_n)$ be a basis of $\mathcal{V}$ and $\gamma$ be a basis of $\mathcal{W}$.\\
 If $S,T\in\LL(\mathcal{V},\mathcal{W})$ and $a\in\mathbb{F}$, then $[aS+T]_{\gamma \leftarrow \beta}=a[S]_{\gamma \leftarrow \beta}+[T]_{\gamma \leftarrow \beta}$.
\begin{proof}
The $j$th column of $[aS+T]_{\gamma \leftarrow \beta}$ is $[(aS+T)\mathbf{v}_j]_{\gamma}=[aS\mathbf{v}_j+T\mathbf{v}_j]_{\gamma}=a[S\mathbf{v}_j]_{\gamma}+[T\mathbf{v}_j]_{\gamma}$, which is the $j$th column of $a[S]_{\gamma \leftarrow \beta}+[T]_{\gamma \leftarrow \beta}$.\\
\hlt{Structural Commentary}: Matrix representation records basis images column-by-column.\\
\hlt{Key Insight}: Linear combinations of maps become entrywise linear combinations of their matrices.
\end{proof}
\end{proposition}

\begin{definition} \hlt{Linear Map Induced by a Matrix}\\
For $A\in\operatorname{Mat}(m,n,\mathbb{F})$, the \hlt{linear map induced by $A$} is the map $T_A:\mathbb{F}^n\to\mathbb{F}^m$ defined by $T_A(\mathbf{x})=A\mathbf{x}$.
\end{definition}

\begin{proposition} \label{prop:matrix-induced-map-linear} \hlt{Matrix-Induced Maps Are Linear}\\
If $A\in\operatorname{Mat}(m,n,\mathbb{F})$, then $T_A:\mathbb{F}^n\to\mathbb{F}^m$ is linear.
\begin{proof}
For $\mathbf{x},\mathbf{y}\in\mathbb{F}^n$ and $a\in\mathbb{F}$, $T_A(a\mathbf{x}+\mathbf{y})=A(a\mathbf{x}+\mathbf{y})=aA\mathbf{x}+A\mathbf{y}=aT_A(\mathbf{x})+T_A(\mathbf{y})$.\\
\hlt{Structural Commentary}: Matrix multiplication preserves addition and scalar multiplication.\\
\hlt{Key Insight}: Every matrix defines a canonical linear map between coordinate spaces.
\end{proof}
\end{proposition}

\begin{corollary} \label{cor:matrix-rank-nullity} \hlt{Matrix Rank-Nullity}\\
Let $A\in\operatorname{Mat}(m,n,\mathbb{F})$. Then $\dim(\ker(A))+\dim(\operatorname{col}(A))=n$ and $\operatorname{nullity}(A)+\operatorname{rank}(A)=n$.\\
If $m=n$, then $\ker(A)=\{\mathbf{0}\}$ if and only if $\operatorname{col}(A)=\mathbb{F}^n$.
\begin{proof}
Apply Theorem \ref{thm:rank-nullity} to the induced linear map $T_A:\mathbb{F}^n\to\mathbb{F}^m$. By definition, $\ker(T_A)=\ker(A)$ and $\operatorname{range}(T_A)=\operatorname{col}(A)$, giving the displayed identities. When $m=n$, the remaining equivalence follows as:
\begin{enumerate}[label=\roman*.]
\setlength{\itemsep}{0pt}
\item[$(\Rightarrow)$] If $\ker(A)=\{\mathbf{0}\}$, then $\dim(\operatorname{col}(A))=n$. Since $\operatorname{col}(A)\subseteq\mathbb{F}^n$, Proposition \ref{prop:dimension-bounds-and-basis-criteria} gives $\operatorname{col}(A)=\mathbb{F}^n$.
\item[$(\Leftarrow)$] If $\operatorname{col}(A)=\mathbb{F}^n$, then $\dim(\operatorname{col}(A))=n$, so the identity gives $\dim(\ker(A))=0$. Hence $\ker(A)=\{\mathbf{0}\}$.
\end{enumerate}
\hlt{Structural Commentary}: Matrix rank-nullity is ordinary rank-nullity written in coordinates.\\
\hlt{Key Insight}: Kernel loss plus column-space dimension equals the number of columns.
\end{proof}
\end{corollary}

\begin{theorem} \label{thm:row-rank-equals-column-rank} \hlt{Row Rank Equals Column Rank}\\
For $A\in\operatorname{Mat}(m,n,\mathbb{F})$, $\dim(\operatorname{row}(A))=\dim(\operatorname{col}(A))$. Equivalently, $\operatorname{rank}(A^t)=\operatorname{rank}(A)$.
\begin{proof}
Let $r=\operatorname{rank}(A)$. If $r=0$, then $\operatorname{col}(A)=\{\mathbf{0}\}$, so every column of $A$ is zero and $A=0$; the result follows. Suppose $r>0$. Choose columns $\mathbf{a}_{j_1},\ldots,\mathbf{a}_{j_r}$ forming a basis of $\operatorname{col}(A)$, and set $B=[\mathbf{a}_{j_1}\ \cdots\ \mathbf{a}_{j_r}]\in\operatorname{Mat}(m,r,\mathbb{F})$.\\
For each column $\mathbf{a}_k$ of $A$, there is a unique $\mathbf{x}_k\in\mathbb{F}^r$ such that $\mathbf{a}_k=B\mathbf{x}_k$. With $X=[\mathbf{x}_1\ \cdots\ \mathbf{x}_n]\in\operatorname{Mat}(r,n,\mathbb{F})$, we have $A=BX$. Therefore $A^t=X^tB^t$, so every column of $A^t$ is a linear combination of the $r$ columns of $X^t$. Hence $\operatorname{rank}(A^t)\leq r=\operatorname{rank}(A)$. Applying the same argument to $A^t$ gives $\operatorname{rank}(A)\leq\operatorname{rank}(A^t)$.\\
\hlt{Structural Commentary}: Column rank controls row rank by factoring $A$ through a basis of its column space.\\
\hlt{Key Insight}: Rows and columns measure the same intrinsic rank, viewed from transposed coordinates.
\end{proof}
\end{theorem}

\begin{perspective} \label{per:row-column-rank-different-hosts} \hlt{Row and Column Spaces Have Different Hosts}\\
For $A\in\operatorname{Mat}(m,n,\mathbb{F})$, $\operatorname{row}(A)\subseteq\mathbb{F}^n$ and $\operatorname{col}(A)\subseteq\mathbb{F}^m$.\\
Hence, when $m\neq n$, Theorem \ref{thm:row-rank-equals-column-rank} asserts only $\dim\operatorname{row}(A)=\dim\operatorname{col}(A)$, not equality of subspaces.
\end{perspective}

\begin{theorem} \label{thm:matrix-four-fundamental-subspaces-dimensions} \hlt{Four Fundamental Matrix Subspaces}\\
Let $A\in\operatorname{Mat}(m,n,\mathbb{F})$ and let $r=\operatorname{rank}(A)$. Then $\operatorname{row}(A),\ker(A) \subseteq \mathbb{F}^n$; $\operatorname{col}(A),\ker(A^t) \subseteq \mathbb{F}^m$,\\
and $\dim(\operatorname{col}(A))=\dim(\operatorname{row}(A))=r$; $\dim(\ker(A))=n-r$; $\dim(\ker(A^t))=m-r$.
\begin{proof}
The inclusions follow from the ambient dimensions of the rows and columns, and from the domains of $A$ and $A^t$. By definition, $\dim(\operatorname{col}(A))=r$; by Theorem \ref{thm:row-rank-equals-column-rank}, $\dim(\operatorname{row}(A))=r$.\\
Rank-nullity gives $\dim(\ker(A))=n-\operatorname{rank}(A)=n-r$. Since $\operatorname{rank}(A^t)=\dim(\operatorname{col}(A^t))=\dim(\operatorname{row}(A))=r$, rank-nullity for $A^t:\mathbb{F}^m\to\mathbb{F}^n$ gives $\dim(\ker(A^t))=m-\operatorname{rank}(A^t)=m-r$.\\
\hlt{Structural Commentary}: One rank determines all four fundamental dimensions.\\
\hlt{Key Insight}: Rank counts what survives; nullity counts what is killed.
\end{proof}
\end{theorem}

\begin{figure}[H]
  \centering
  \scalebox{0.9}{\begin{tikzpicture}[
  scale=0.90,
  transform shape,
  >=Stealth,
  every node/.style={font=\small},
  host/.style={draw=gray!50, fill=gray!3, rounded corners=3pt, line width=0.55pt,
    minimum width=4.05cm, minimum height=2.05cm},
  entry/.style={draw=gray!60, rounded corners=2pt, line width=0.5pt,
    minimum width=1.72cm, minimum height=0.78cm, align=center},
  active/.style={entry, fill=teal!8},
  null/.style={entry, fill=blue!7},
  map/.style={->, gray!65!black, line width=0.75pt}
]
  \node[host] (Fn) at (0,0) {};
  \node[anchor=north west] at (-1.82,0.88) {$\mathbb{F}^n$};
  \node[host] (Fm) at (5.6,0) {};
  \node[anchor=north west] at (3.78,0.88) {$\mathbb{F}^m$};

  \node[null] at (-0.98,-0.05) {$\ker(A)$\\[-1pt]{\footnotesize $n-r$}};
  \node[active] at (0.98,-0.05) {$\operatorname{row}(A)$\\[-1pt]{\footnotesize $r$}};
  \node[active] at (4.62,-0.05) {$\operatorname{col}(A)$\\[-1pt]{\footnotesize $r$}};
  \node[null] at (6.58,-0.05) {$\ker(A^t)$\\[-1pt]{\footnotesize $m-r$}};

  \draw[map] (2.1,0.4) -- (3.5,0.4) node[midway, above] {$A$};
  \draw[map] (3.5,-0.4) -- (2.1,-0.4) node[midway, below] {$A^t$};
\end{tikzpicture}
}
  \caption{Four fundamental subspaces by host coordinate space.}
  \label{fig:four-fundamental-subspaces-hosts}
\end{figure}

\begin{corollary} \label{cor:homogeneous-system-more-variables-than-equations} \hlt{Homogeneous Systems with More Variables Than Equations}\\
If $A\in\operatorname{Mat}(m,n,\mathbb{F})$ with $n>m$, then the homogeneous system $A\mathbf{x}=\mathbf{0}$ has a nonzero solution.
\begin{proof}
The induced map $T_A:\mathbb{F}^n\to\mathbb{F}^m$ has domain dimension larger than codomain dimension, so Corollary \ref{cor:map-to-smaller-space-not-injective} shows that $T_A$ is not injective.\\
Hence $\ker(A)=\ker(T_A)$ contains a nonzero vector $\mathbf{x}$, and this vector satisfies $A\mathbf{x}=\mathbf{0}$.\\
\hlt{Structural Commentary}: Extra variables force a nontrivial kernel.\\
\hlt{Key Insight}: A nontrivial kernel is exactly a nonzero homogeneous solution.
\end{proof}
\end{corollary}

\begin{corollary} \label{cor:inhomogeneous-system-more-equations-than-variables} \hlt{Inhomogeneous Systems with More Equations Than Variables}\\
If $A\in\operatorname{Mat}(m,n,\mathbb{F})$ with $m>n$, then there exists $\mathbf{b}\in\mathbb{F}^m$ such that $A\mathbf{x}=\mathbf{b}$ has no solution.
\begin{proof}
The induced map $T_A:\mathbb{F}^n\to\mathbb{F}^m$ has domain dimension smaller than codomain dimension, so Corollary \ref{cor:map-to-larger-space-not-surjective} shows that $T_A$ is not surjective. Choose $\mathbf{b}\in\mathbb{F}^m\setminus\operatorname{range}(T_A)$. Then no $\mathbf{x}\in\mathbb{F}^n$ satisfies $A\mathbf{x}=\mathbf{b}$.\\
\hlt{Structural Commentary}: Too few variables cannot reach every possible right-hand side.\\
\hlt{Key Insight}: A missing range vector becomes an inconsistent inhomogeneous system.
\end{proof}
\end{corollary}

\begin{proposition} \label{prop:matrix-representation-action} \hlt{Matrix Representation Acts on Coordinates}\\
Let $\beta=(\mathbf{v}_1,\ldots,\mathbf{v}_n)$ be a basis of $\mathcal{V}$, let $\gamma$ be a basis of $\mathcal{W}$, and let $T\in\LL(\mathcal{V},\mathcal{W})$.\\
Then for every $\mathbf{v} \in \mathcal{V}$, $[T(\mathbf{v})]_{\gamma} = [T]_{\gamma \leftarrow \beta}[\mathbf{v}]_{\beta}$.
\begin{proof}
Write $\mathbf{v} = a_1\mathbf{v}_1 + \cdots + a_n\mathbf{v}_n$. Then linearity gives $T(\mathbf{v}) = a_1T(\mathbf{v}_1) + \cdots + a_nT(\mathbf{v}_n)$ and the right-hand side is exactly the linear combination of the columns of $[T]_{\gamma \leftarrow \beta}$ with coefficients $a_1,\ldots,a_n$.\\
\hlt{Structural Commentary}: Matrix representation turns applying $T$ into multiplying coordinate columns.\\
\hlt{Key Insight}: The columns of $[T]_{\gamma \leftarrow \beta}$ are the coordinate images of the basis vectors.
\end{proof}
\end{proposition}

\begin{theorem} \label{thm:invertible-matrix-columns-basis} \hlt{Columns of an Invertible Matrix Form a Basis}\\
Let $A=[\mathbf{a}_1\ \cdots\ \mathbf{a}_n]\in\operatorname{Mat}(n,n,\mathbb{F})$ be invertible. Then $(\mathbf{a}_1,\ldots,\mathbf{a}_n)$ is a basis of $\mathbb{F}^n$.
\begin{proof}
The columns span $\mathbb{F}^n$ because for any $\mathbf{y}\in\mathbb{F}^n$, writing $\mathbf{x}=A^{-1}\mathbf{y}$ gives $\mathbf{y}=A\mathbf{x}=x_1\mathbf{a}_1+\cdots+x_n\mathbf{a}_n$.\\
They are linearly independent because if $c_1\mathbf{a}_1+\cdots+c_n\mathbf{a}_n=\mathbf{0}$, then $A\mathbf{c}=\mathbf{0}$, so $\mathbf{c}=A^{-1}\mathbf{0}=\mathbf{0}$.\\
\hlt{Structural Commentary}: An invertible matrix is a coordinate isomorphism of $\mathbb{F}^n$. Its columns are the images of the standard coordinate directions.\\
\hlt{Key Insight}: Invertibility forces the columns to span and prevents nontrivial column relations.
\end{proof}
\end{theorem}

\begin{theorem} \label{thm:square-one-sided-inverse-is-two-sided} \hlt{One-Sided Inverses of Square Matrices}\\
Let $A,B\in\operatorname{Mat}(n,n,\mathbb{F})$. If $AB=I_n$, then $BA=I_n$.
\begin{proof}
If $B\mathbf{x}=\mathbf{0}$, then $\mathbf{x}=AB\mathbf{x}=\mathbf{0}$, so the induced map $T_B:\mathbb{F}^n\to\mathbb{F}^n$ is injective. By Theorem \ref{thm:finite-dimensional-operator-invertible-equivalences}, $T_B$ is surjective. For any $\mathbf{y}\in\mathbb{F}^n$, choose $\mathbf{x}$ with $\mathbf{y}=B\mathbf{x}$. Then $BA\mathbf{y}=BAB\mathbf{x}=B\mathbf{x}=\mathbf{y}$.\\
Thus $BA$ acts as the identity on every vector of $\mathbb{F}^n$, so $BA=I_n$.\\
\hlt{Structural Commentary}: In finite dimensions, injectivity and surjectivity coincide for square linear maps. A left inverse makes $B$ injective, hence onto, which forces the missing inverse relation.\\
\hlt{Key Insight}: A one-sided inverse of a square matrix is automatically two-sided.
\end{proof}
\end{theorem}

\begin{proposition} \label{prop:composition-and-matrix-multiplication} \hlt{Composition and Matrix Multiplication}\\
If $T \in \LL(\mathcal{U}, \mathcal{V})$ and $S \in \LL(\mathcal{V}, \mathcal{W})$, then $[ST]_{\delta \leftarrow \alpha} = [S]_{\delta \leftarrow \beta}[T]_{\beta \leftarrow \alpha}$ for any bases $\alpha$ of $\mathcal{U}$, $\beta$ of $\mathcal{V}$, and $\delta$ of $\mathcal{W}$.
\begin{proof}
For every $\mathbf{u}\in\mathcal{U}$, $[(ST)\mathbf{u}]_{\delta}=[S]_{\delta \leftarrow \beta}[T\mathbf{u}]_{\beta}=[S]_{\delta \leftarrow \beta}[T]_{\beta \leftarrow \alpha}[\mathbf{u}]_{\alpha}$.\\
Since coordinate columns $[\mathbf{u}]_{\alpha}$ range over all of $\mathbb{F}^{\dim(\mathcal{U})}$, the representing matrices are equal.\\
\hlt{Structural Commentary}: Matrix multiplication is the coordinate shadow of composition of linear maps.\\
\hlt{Key Insight}: Applying $T$ then $S$ becomes multiplying first by the matrix of $T$, then by the matrix of $S$.
\end{proof}
\end{proposition}

\begin{theorem} \label{thm:linear-maps-and-matrices-same-data} \hlt{Linear Maps and Matrices Are the Same Data}\\
If $\dim(\mathcal{V}) = n$ and $\dim(\mathcal{W}) = m$, then choosing bases of $\mathcal{V}$ and $\mathcal{W}$ gives isomorphism $\LL(\mathcal{V}, \mathcal{W}) \cong \operatorname{Mat}(m,n,\mathbb{F})$.\\
Consequently, $\dim(\LL(\mathcal{V}, \mathcal{W})) = mn$.
\begin{proof}
Fix bases $\beta$ of $\mathcal{V}$ and $\gamma$ of $\mathcal{W}$. The map $T\mapsto [T]_{\gamma \leftarrow \beta}$ is linear. It is injective because a linear map whose matrix is zero sends every basis vector of $\mathcal{V}$ to $\mathbf{0}$, hence is the zero map. It is surjective because, for any $A=[a_{ij}]\in\operatorname{Mat}(m,n,\mathbb{F})$, there is a unique linear map satisfying $T(\mathbf{v}_j)=a_{1j}\mathbf{w}_1+\cdots+a_{mj}\mathbf{w}_m$, $j=1,\ldots,n$.\\
Thus $\LL(\mathcal{V},\mathcal{W})\cong\operatorname{Mat}(m,n,\mathbb{F})$, whose standard entry matrices form a basis of length $mn$.\\
\hlt{Structural Commentary}: Choosing bases converts every linear map into one and only one array of scalars.\\
\hlt{Key Insight}: A matrix is exactly the list of coordinate images of the domain basis vectors.
\end{proof}
\end{theorem}

\begin{motivation} \hlt{Coordinate Encoding}\\
This is the key bridge between abstraction and computation. The linear map is basis-free; the matrix is what the map looks like after coordinates have been chosen.
\end{motivation}
